% =================================================================================================
% CS 401R final report (spec §9). NeurIPS 2025 style (official neurips_2025.sty), single column.
%
% Build:  make report            (needs `brew install tectonic`, or latexmk from MacTeX)
%   or:   Overleaf: upload this folder plus figures/ and tables/, then set \figdir and \tabdir below.
%
% PAGE BUDGET (instructor, Nov 19). 4-6 pages; figures count, references and the appendix don't. The limit is soft,
% but every paragraph has to earn its place and figures must not be blown up to fill space.
%   Target: ~2,600 words of prose + ~1.3 pages of figures and tables = about 5.5 pages before the
%   references. A page of prose in this style holds about 700-750 words (measured on the first build).
%   Figures: full text width, 1.6-1.7 in tall (the height passed to \figfile is a maximum). Each one
%   is referenced in the text and makes one point; one that doesn't change the argument goes to the
%   appendix. The appendix holds what a careful reader might check (hyperparameters, validation,
%   secondary figures), not overflow that the main text's argument depends on.
%
% PLACEHOLDERS. \filler{N} draws N words of gray stand-in text, so the compiled draft shows its real
% length before you've written anything. Delete a section's \filler as you write that section.
% Red \draft{...} notes say what goes where. \notesfalse hides just the notes (the length you'd hand
% in if every section hit its budget); \placeholdersfalse hides notes and filler, leaving only real text.
% Figures and tables that scripts/analyze.py hasn't written yet appear as boxes of their final size.
% =================================================================================================

\documentclass{article}

% `preprint` shows the author block and turns off the anonymous-submission line numbers.
\usepackage[preprint]{neurips_2025}

\usepackage[utf8]{inputenc}
\usepackage[T1]{fontenc}
\usepackage{url}
\usepackage{booktabs}
\usepackage{amsmath}
\usepackage{amsfonts}
\usepackage{nicefrac}
\usepackage{microtype}
\usepackage{xcolor}
\usepackage{graphicx}
\usepackage{hyperref}  % load last
\hypersetup{colorlinks=true, allcolors=blue!50!black}

% The style's first-page footer reads "Preprint." This is a class report, so say that instead.
\makeatletter
\renewcommand{\@noticestring}{CS 401R: Reinforcement Learning. Final project report, December 2026.}
\makeatother

% --- condition names (rename here once and the whole paper follows) -------------------------------
\newcommand{\envT}{\textsf{env}}          % opponent-as-environment targets
\newcommand{\negT}{\textsf{negamax}}      % negamax targets
\newcommand{\naiveS}{\textsf{naive}}      % latest vs latest
\newcommand{\poolS}{\textsf{pool}}        % past checkpoints

% --- placeholders (see header) ---------------------------------------------------------------------
\newif\ifplaceholders \placeholderstrue
\newif\ifnotes \notestrue

\newcommand{\draft}[1]{\ifplaceholders\ifnotes{\color{red!60!black}\sffamily\small[#1]}\par\fi\fi}

\newcount\fillerwords \newcount\fillerpara
\newcommand{\filler}[1]{%
  \ifplaceholders\par{\color{black!30}\fillerwords=0\relax\fillerpara=0\relax
    \loop\ifnum\fillerwords<#1\relax
      lorem ipsum dolor sit amet, consectetur adipiscing elit sed do %
      \advance\fillerwords by 10\relax
      \advance\fillerpara by 10\relax
      \ifnum\fillerpara>119 \par\fillerpara=0\relax\fi
    \repeat\par}\fi}

% --- figures and tables written by scripts/analyze.py ----------------------------------------------
% File names are the contract with analyze.py (see its docstrings). On Overleaf, set these to
% `figures` and `tables` after uploading those folders next to main.tex.
\newcommand{\figdir}{../figures}
\newcommand{\tabdir}{../tables}

\newcommand{\missingbox}[2]{% {height}{label}
  \fbox{\parbox[c][#1][c]{\dimexpr\linewidth-2\fboxsep-2\fboxrule\relax}{\centering\small\color{black!50}%
    \texttt{\detokenize{#2}}\\[2pt]not generated yet}}}

% \figfile{<exp>/<file>.pdf}{height}: the figure if it exists, else a box of that height.
\newcommand{\figfile}[2]{%
  \IfFileExists{\figdir/#1}%
    {\includegraphics[width=\linewidth,height=#2,keepaspectratio]{\figdir/#1}}%
    {\missingbox{#2}{figures/#1}}}

% \tabfile{<exp>/<file>.tex}{height}: \input the generated tabular if it exists, else a box.
\newcommand{\tabfile}[2]{%
  \IfFileExists{\tabdir/#1}{\input{\tabdir/#1}}{\missingbox{#2}{tables/#1}}}

% =================================================================================================

\title{Self-Play with Deep Q-Learning in Othello}

\author{%
  Johnny Goodlad\\
  CS 401R: Reinforcement Learning\\
  % TODO: affiliation and email, if the course wants them
}

\begin{document}

\maketitle

% ~150 words
\begin{abstract}
  \draft{Abstract (150 words): the question, the $2\times2$ design, the main result for each
  hypothesis, and the one-line takeaway about self-play win rate.}
  \filler{150}
\end{abstract}

% =================================================================================================
% ~350 words. Spec §1 and the proposal's "Why this question matters".
\section{Introduction}

\draft{Introduction (350 words): the question; why self-play win rate is a weak progress signal
(it hovers near 50\% whether the agent improves, stalls, or forgets); the $2\times2$ design;
hypotheses H1--H4 as a short contributions list; point to Figure~\ref{fig:headline}.}
\filler{175}

\begin{figure}[t]
  \centering
  \figfile{e2-6x6/f3_selfplay_vs_external.pdf}{1.7in}
  \caption[F3]{Self-play score (current network vs.\ previous checkpoint) and score against the
  positional player over training on $6\times6$. TODO: one-sentence takeaway.}
  \label{fig:headline}
\end{figure}

\filler{175}

% =================================================================================================
% ~250 words. Spec §10 references; D1's AlphaGo/AlphaZero contrast.
\section{Background and related work}

\draft{Background (250 words): self-play \citep{tesauro1995temporal,silver2016mastering,silver2018general};
opponent pools and leagues \citep{bansal2018emergent,lanctot2017unified,vinyals2019grandmaster,berner2019dota},
including AlphaGo sampling past opponents vs.\ AlphaZero using only the latest network;
DQN and Double DQN \citep{mnih2015human,vanhasselt2016double};
checkpoint regret \citep{samvelyan2024madrid} on TiZero's checkpoints \citep{lin2023tizero};
a source for the $4\times4$ and $6\times6$ solved results (still to confirm).}
\filler{250}

% =================================================================================================
% ~400 words plus the equations. Spec §3 D1 (two views), D3 (passes), §4.5 (targets).
\section{Formulation}

\draft{Formulation (400 words): the MDP from the learner's view (states, actions including pass,
terminal $\pm1/0$ reward, $\gamma=1$). The two views of the opponent: under \envT{} the opponent is
part of the transition function, so $Q$ is a best response to the opponents it trained against;
under \negT{}, $Q$ estimates the minimax value whatever the opponent does. Why the pool should
matter more for \envT{} (H4). Passes: an explicit action with a sign flip under \negT{}, absorbed
into the dynamics under \envT{}.}
\filler{225}

Both target types share Double DQN action selection; they differ only in the sign of the bootstrap
term and in what $s'$ is:
\begin{align}
  a^\star &= \arg\max_{a' \in \mathcal{A}(s')} Q_\phi(s', a'), \label{eq:argmax}\\
  y &= \begin{cases}
    r & s' \text{ terminal},\\
    +\,Q_{\bar\phi}(s', a^\star) & \text{\envT{}: } s' \text{ is the learner's next turn},\\
    -\,Q_{\bar\phi}(s', a^\star) & \text{\negT{}: } s' \text{ is the opponent's turn}.
  \end{cases} \label{eq:target}
\end{align}

\filler{125}

% =================================================================================================
% ~250 words. Spec §4.5-4.6, D2.
\section{Method}

\draft{Method (250 words): Q-network (conv trunk, no batch norm, $\tanh$ output over $N^2+1$
actions); Double DQN with Huber loss; illegal-move masking; a random D4 symmetry at sample time;
only the learner's moves are stored, in all four conditions; the pool sampling rule.
Full hyperparameters are in Appendix~\ref{app:hparams}.}
\filler{250}

% =================================================================================================
% ~350 words. Spec §4.3-4.4, §4.7, D4-D6.
\section{Experimental setup}

\begin{table}[t]
  \caption[Design]{The $2\times2$ design. All other settings are shared across conditions.}
  \label{tab:design}
  \centering
  \begin{tabular}{lcc}
    \toprule
    & \naiveS{} (latest vs.\ latest) & \poolS{} (past checkpoints) \\
    \midrule
    \envT{} targets (opponent as environment) & \envT-\naiveS & \envT-\poolS \\
    \negT{} targets & \negT-\naiveS & \negT-\poolS \\
    \bottomrule
  \end{tabular}
\end{table}

\draft{Setup (350 words): boards and seeds ($4\times4$: 5 seeds; $6\times6$: 3); game and gradient
budgets; reference players (random, greedy, positional); the exact $4\times4$ solver and the
$6\times6$ endgame solver; opening banks played with both colors; score with 95\% Wilson intervals
\citep{wilson1927probable}; self-play score = current vs.\ previous checkpoint; per-run summary
numbers; checkpoint regret (Eq.~\ref{eq:regret}) against self-history and held-out references.}
\filler{225}

For a position $\theta$ and a target agent $T$, a reference policy $R$ has regret
\begin{equation}
  \mathrm{reg}_R(\theta) = V^\theta(R, T) - V^\theta(T, T) \in \{-2, \dots, 2\},
  \label{eq:regret}
\end{equation}
where $V^\theta(X, T)$ is the outcome for the side to move when $X$ plays it from $\theta$ against $T$.

\filler{125}

% =================================================================================================
% ~500 words total, plus the main-text figures and table.
\section{Results}

\subsection{\texorpdfstring{$4\times4$}{4x4}: accuracy against exact play}
\draft{(150 words) Figure~\ref{fig:acc4}; H1 on $4\times4$; Q-vs-solver error for \negT{} only.}
\filler{150}

\begin{figure}[t]
  \centering
  \figfile{e1-4x4/f1_accuracy.pdf}{1.6in}
  \caption[F1]{Optimal-move accuracy on $4\times4$ against games played, for \envT{} (left) and
  \negT{} (right) targets; \naiveS{} vs.\ \poolS{}, bands span the 5 seeds' range, dotted line = chance.
  TODO: takeaway.}
  \label{fig:acc4}
\end{figure}

\subsection{\texorpdfstring{$6\times6$}{6x6}: strength against reference players}
\draft{(175 words) Figure~\ref{fig:scores6} and Table~\ref{tab:summary}; H1 and H2; the
\poolS{}$-$\naiveS{} gap under each target type (H4).}
\filler{175}

\begin{figure}[t]
  \centering
  \figfile{e2-6x6/f2_scores.pdf}{1.6in}
  \caption[F2]{Score against the positional player on $6\times6$, same layout as
  Figure~\ref{fig:acc4}; bands over 3 seeds. TODO: takeaway.}
  \label{fig:scores6}
\end{figure}

\begin{table}[t]
  \caption[T1]{Per-run summary on $6\times6$, mean $\pm$ std over seeds: final score, peak, maximum
  drawdown, forgetting events, Spearman $\rho$ between self-play and external score, and the
  \poolS{}$-$\naiveS{} gap per target type.}
  \label{tab:summary}
  \centering
  \tabfile{e2-6x6/t1_summary.tex}{1.0in}
\end{table}

\subsection{What self-play hides}
\draft{(175 words) Figure~\ref{fig:headline} and Figure~\ref{fig:regret}; H3; positions where earlier
checkpoints beat the final agent; self-history vs.\ held-out regret.}
\filler{175}

\begin{figure}[t]
  \centering
  \figfile{e2-6x6/f7_regret.pdf}{1.7in}
  \caption[F7]{Left: self-history regret by reference checkpoint for the final \naiveS{} and \poolS{}
  agents. Right: regret over (empty squares $\times$ disc differential). TODO: takeaway.}
  \label{fig:regret}
\end{figure}

% =================================================================================================
% ~200 words. Spec D1 note (replay buffer as an implicit pool).
\section{Discussion}

\draft{Discussion (200 words): what self-play metrics hide; when pools matter (the H4 result); the
replay buffer as an implicit short-window pool under \envT-\naiveS{}; anything from the presentation
Q\&A.}
\filler{200}

% =================================================================================================
% ~100 words. Can merge into Discussion if space is tight.
\section{Limitations and future work}

\draft{Limitations (100 words): small boards, few seeds, one seed in the demo; future work: $8\times8$,
prioritized pool sampling, search, MAP-Elites over positions.}
\filler{100}

% =================================================================================================
% ~75 words.
\section{Conclusion}

\draft{Conclusion (75 words): answer the question in two or three sentences.}
\filler{75}

% =================================================================================================

{\small
\bibliographystyle{plainnat}
\bibliography{references}
}

% =================================================================================================
\appendix

\section{Hyperparameters}
\label{app:hparams}

% From spec §4.5. Update to the values actually run (the 6x6 game budget is fixed on Nov 24).
\begin{table}[h]
  \caption[Hyperparameters]{Hyperparameters, shared by all four conditions.}
  \label{tab:hparams}
  \centering
  \small
  \begin{tabular}{lll}
    \toprule
    & $4\times4$ & $6\times6$ \\
    \midrule
    Network trunk & 3 conv layers, 64 channels & 4 conv layers, 64 channels \\
    Training games & 30k & $\sim$200k \\
    Parallel games & 128 & 256 \\
    Replay capacity (transitions) & 50k & 300k \\
    Evaluation and checkpoint interval & 1k games & 5k games \\
    Seeds & 5 & 3 \\
    \midrule
    Batch size, optimizer & \multicolumn{2}{l}{256, Adam with learning rate $3\times10^{-4}$} \\
    Gradient steps & \multicolumn{2}{l}{1 per 16 new learner transitions} \\
    Target network & \multicolumn{2}{l}{hard update every 1{,}000 gradient steps} \\
    Exploration & \multicolumn{2}{l}{$\varepsilon$ from 1.0 to 0.05 over the first 20\% of games} \\
    Loss & \multicolumn{2}{l}{Huber, $\gamma = 1$, gradient-norm clip 10} \\
    Pool sampling & \multicolumn{2}{l}{latest with $p = 0.5$, else a uniform past checkpoint} \\
    \bottomrule
  \end{tabular}
\end{table}

\section{Validation}
\label{app:validation}

\draft{Engine: 8$\times$8 perft through depth 6 and 10k random games; TypeScript engine and ONNX parity
tests; the $4\times4$ start-position value against the published result; endgame-solver checks.}
\filler{150}

\section{Additional results}
\label{app:extra}

\draft{Cross-play heatmaps (F4), $6\times6$ endgame accuracy (F5), checkpoint vs.\ true regret on
$4\times4$ (F8), the $4\times4$ summary table. F6 (Q heatmaps on teaching positions) has no
analyze.py function yet; add one or drop it.}

\begin{figure}[h]
  \centering
  \figfile{e2-6x6/f4_crossplay.pdf}{2.0in}
  \caption[F4]{Cross-play between 15 checkpoints per run, \naiveS{} vs.\ \poolS{}, $6\times6$.}
  \label{fig:crossplay}
\end{figure}

\begin{figure}[h]
  \centering
  \figfile{e2-6x6/f5_endgame.pdf}{1.8in}
  \caption[F5]{Endgame accuracy on 500 solved $6\times6$ positions against games played.}
  \label{fig:endgame}
\end{figure}

\begin{figure}[h]
  \centering
  \figfile{e1-4x4/f8_true_vs_checkpoint_regret.pdf}{1.8in}
  \caption[F8]{Checkpoint-estimated regret vs.\ true regret against the optimal policy, $4\times4$.}
  \label{fig:trueregret}
\end{figure}

\begin{table}[h]
  \caption[T1 4x4]{Per-run summary on $4\times4$, same columns as Table~\ref{tab:summary}.}
  \label{tab:summary4}
  \centering
  \tabfile{e1-4x4/t1_summary.tex}{1.2in}
\end{table}

\section{Demo}
\label{app:demo}

The trained agents can be played in the browser at \url{https://othello.johnnygoodlad.dev}.
% Spec §5.3.
\draft{Which condition and seed the demo uses, and that it is one seed.}

\end{document}
